\section{Coding Theory}\label{sec:coding_theory}

This section contains definitions and theorems about coding theory as they are used in the rest of the library.

\begin{definition}[Code Distance]
    \label{def:code_distance}
    \lean{Code.dist}
    \leanok
    The distance of a code $\code\subseteq\Sigma^{\evaldomain}$ is the least Hamming distance
    between two distinct codewords; it is $0$ if $\code$ has fewer than two elements.
\end{definition}

\begin{definition}[Distance from a Code]
    \label{def:distance_from_code}
    \lean{Code.distFromCode}
    \leanok
    The distance $\Delta(u,\code)$ of a word $u$ from a code $\code$ is the least Hamming distance
    between $u$ and a codeword, valued in $\mathbb{N}\cup\{\infty\}$ (it is $\infty$ when
    $\code$ is empty).
\end{definition}

\begin{definition}[Generator Matrix]
    \label{def:generator_matrix}
    \lean{LinearCode.fromRowGenMat, LinearCode.fromColGenMat}
    \leanok
    A matrix $G$ generates the linear code spanned by its rows (equivalently, the image of
    $v\mapsto vG$); the column variant takes the image of $v\mapsto Gv$.
\end{definition}

\begin{definition}[Parity Check Matrix]
    \label{def:parity_check_matrix}
    \lean{LinearCode.byCheckMatrix}
    \leanok
    A matrix $H$ defines the linear code $\{c : Hc = 0\}$, the kernel of $c\mapsto Hc$.
\end{definition}

\begin{definition}[Code]
    \label{def:code}
    A code of length $\evaldomain$ over an alphabet $\Sigma$ is a set
    $\code\subseteq\Sigma^{\evaldomain}$ of words. The formalization uses
    $\mathsf{Set}\,(\evaldomain\to\Sigma)$ directly and has no separate declaration for it.
\end{definition}

\begin{definition}[Linear Code]
    \label{def:linear_code}
    \lean{LinearCode}
    \leanok
    A linear code of length $\evaldomain$ over $\field$ is an $\field$-submodule of
    $\field^{\evaldomain}$.
\end{definition}

\begin{definition}[Interleaved Code]
    \label{def:interleaved_code}
    \lean{Code.interleavedCodeSet}
    \leanok
    The $\kappa$-fold interleaving of a code $\code$ is the set of
    $\evaldomain\times\kappa$ matrices each of whose columns lies in $\code$.
\end{definition}

\begin{definition}[Reed-Solomon Code]
    \label{def:reed_solomon_code}
    \lean{ReedSolomon.code}
    \leanok
    For an injective map $\evaldomain\hookrightarrow\field$ and $\degree\in\mathbb{N}$,
    $\rscode[\field,\evaldomain,\degree]$ is the image of the polynomials of degree less than
    $\degree$ under evaluation on $\evaldomain$.
\end{definition}

\begin{definition}[Smooth Reed-Solomon Code]
    \label{def:smooth_rs_code}
    \lean{ReedSolomon.code, Domain.SmoothCosetFftDomain}
    \leanok
    A smooth Reed--Solomon code is a Reed--Solomon code whose evaluation domain is a smooth coset
    FFT domain: a coset of a multiplicative subgroup of $\field^\times$ of order $2^n$.
\end{definition}

\begin{definition}[Constrained Code]
    \label{def:constrained_code}
    \lean{ReedSolomon.constrainedCode}
    \leanok
    For a smooth domain $\evaldomain$ of size $2^n$, $m\in\mathbb{N}$, a weight polynomial
    $\hat{w}\in\field[Z,X_1,\ldots,X_m]$ and $\sigma\in\field$, the constrained code
    $\crscode[\field,\evaldomain,m,\hat{w},\sigma]$ consists of the codewords
    $f\in\rscode[\field,\evaldomain,2^m]$ whose multilinear polynomial $\hat{f}$ satisfies
    $\sum_{b\in\{0,1\}^m}\hat{w}(\hat{f}(b),b)=\sigma$.
\end{definition}

\begin{definition}[Multi-constrained Code]
    \label{def:multi_constrained_code}
    \lean{ReedSolomon.multiConstrainedCode}
    \leanok
    The multi-constrained code for weights $\hat{w}_0,\ldots,\hat{w}_{t-1}$ and targets
    $\sigma_0,\ldots,\sigma_{t-1}$ consists of the codewords of $\rscode[\field,\evaldomain,2^m]$
    that satisfy every one of the $t$ weight constraints.
\end{definition}

\begin{definition}[Proximity Measure]
    \label{def:proximity_measure}
    \lean{ProximityGap.proximityMeasure}
    \leanok
    \uses{def:distance_from_code}
    For words $u,v$, a code $\code$ and $d\in\mathbb{N}$, the proximity measure is the number of
    $r\in\field$ with $\Delta(r u+(1-r)v,\code)\le d$.
\end{definition}

\begin{definition}[Proximity Gap]
    \label{def:proximity_gap}
    \lean{ProximityGap.proximityGap}
    \leanok
    \uses{def:distance_from_code}
    A code $\code$ has a proximity gap at distance $d$ with bound $b$ if, for all words $u,v$
    whose proximity measure exceeds $b$, the pair $(u,v)$ is within distance $d$ of the
    2-fold interleaved code of $\code$.
\end{definition}

\begin{definition}[List Decodability]
    \label{def:list_decodable}
    \lean{Code.IsListDecodable}
    \leanok
    \uses{def:distance_from_code,def:code}
    A code $\code$ is $(r,\ell)$-list decodable if every word has at most $\lfloor\ell\rfloor$
    codewords within relative Hamming distance $r$.
\end{definition}

\begin{definition}[List of Close Codewords]
    \label{def:list_close_codewords}
    \lean{Code.closeCodewordsRel}
    \leanok
    \uses{def:distance_from_code,def:code}
    For a word $y$ and $r\in\mathbb{R}$, $\Lambda(\code,y,r)$ is the set of codewords within
    relative Hamming distance $r$ of $y$.
\end{definition}

\subsection{Theorems}

\begin{theorem}[List decoding of Reed--Solomon codes up to the Johnson bound]
    \label{thm:rs_list_decoding_johnson}
    \lean{ReedSolomon.listDecodable_reedSolomon}
    \leanok
    \uses{def:reed_solomon_code,def:list_decodable}
    Let $\evaldomain$ be a nonempty finite domain embedded in a field $\field$, let $\degree>0$,
    and let $\rate=\min\{\degree,|\evaldomain|\}/|\evaldomain|$ be the rate of
    $\rscode[\field,\evaldomain,\degree]$. For every real $\eta>0$, the code
    $\rscode[\field,\evaldomain,\degree]$ is
    $\bigl(1-\sqrt{\rate}-\eta,\ 1/(2\eta\sqrt{\rate})\bigr)$-list decodable. The field need not
    be finite, and the bound does not depend on $|\field|$.
\end{theorem}
\begin{proof}
    \leanok
    Two distinct codewords agree in fewer than $\degree$ positions, so their agreement is at
    most $\min\{\degree,|\evaldomain|\}=\rate|\evaldomain|$. A counting bound for any code whose
    pairwise agreements are this small gives at most $1/(2\eta\sqrt{\rate})$ codewords within
    relative distance $1-\sqrt{\rate}-\eta$ of any word. This is Theorem~4.3 of~\cite{ACFY24}.
\end{proof}

The proximity-gap and mutual correlated agreement results for Reed--Solomon codes are stated
where the protocols use them: Theorem~\ref{thm:proximity_gap} for random linear combinations,
and Lemma~\ref{lemma:mca_rs_udr} and Theorems~\ref{thm:mca_rs_johnson},
\ref{thm:mca_rs_capacity} and~\ref{thm:power_batching_capacity} for mutual correlated agreement
in the unique-decoding, Johnson and capacity regimes.
